% Nouveau EGA V (cote 25): archive cote 25, site page 54 (PDF page 55).
% typescript. Draft compared: none.
% AI-transcribed by Claude Opus 5.5 (claude-opus-5-5, adaptive thinking, effort high), Message Batches API, batch ega25b (msgbatch_01A7n5YuXubUPNK1zsepi8RG), 2026-11-13; original AI transcription, subsequently checked against both supplied scan parts; see PAGES_CHECKED.tsv and CHECK_DECISIONS.tsv.
\scanpage{25/54}
\mspage{\hand{49}}

dans la version présente  avec L\unsure{,|.} d'avoir une condition néc et suffisan\xout{t}%
te \unsure{d|p}'irréductibilité géom. \hand{de $Y_\xi$, $\xi$ gén. dans $L$.}

Du point de vue cunutesque, \xout{il} et en termes de théorie des
corps, il y a lieu d'exprimer la condition \xout{(suffisante)} en termes de
degré de transcendance de la f\unsure{\c{a}|a}çon suivante: on choisit un ``hyperplan
à l'infini'' ne contenant ni $L^{o}$ ni X, et on se place dans son \unsure{c|v}omplé%
mentaire i.e. sur \unsure{l|h}e schéma affine  type essentiellement. On choisit
une base de l'\unsure{éspace|espace} des formes linéaires nulles sur $L^{o}$ , \xout{savoir}
\xout{et on regarde} soit $T_{1},\ldots,T_{p}$ ( p = codim ($L^{o}$,P) ), et on considère
leurs images inverses $t_{1},..,t_{p}$ dans le corps des fractions K de X
(\inserted{X} supposé intègre). Un au moins des $t_{i}$, mettons $t_{1}$ , est $\hand{\neq} 0$, considé-%
rons alors $a_{1} = t_{2}/\unsure{t|l}_{1}$ ,..., $a_{p-1} = t_{p}/t_{1}$ . Alors dim g(X') n'est
autre que le degré de transcendance de $k(a_{1},..,a_{p-1}) \hand{\subset} K$ sur k. Donc
si ce degré de transcendance est $\hand{\geqslant} 2$, O.K. \xout{De même s'il} S'il est 1,
alors il faut exiger que sur $\overset{\hand{\rule{0.6em}{0.35pt}}}{k}$ , f(X) soit contenu dans une sous-varié%
té linéaire de P contenant $L^{o}$ et de dimension une de plus\hand{, et que la
fibre générique de $g : X' \to \overline{g(X')}$ soit géom. irréd.}

Supposons que $L^{o}$ soit de dimension q, alors les fibres de u:
$P - L^{o} \hand{\to} Q$ sont de dimmension q+1, \inserted{donc celles de $g$ sont de dim $\leq q+1$,} et par suite on a
\[ \dim g(X') \hand{\geqslant} \dim f(X) - (q+1) , \]
donc la condition dimensionnelle sur g(X') est vérifiée pourvu que
\[ \dim f(X) \hand{\geqslant} q+\hand{3} \quad . \]
Lorsque q=0, on trouve le fait signalé dans 8.17. a). \inserted[typed]{\xout{du}} \xout{Lorsque q=0,}

Revenant aux conditions de 8.18. on voit que la condition a)
relative à un $x \hand{\in} T$ est sat\unsure{i|m}sfaite pourvu que X ne soit pas ``conique''
en x relativement à \unsure{$f$|£}'', dans un sens évident.

A voir si ces derniers développements Bertiniesques ne viendrai-%
ent pas mieux au n\textsuperscript{o} suivant du changement d'immersion projective. \margin{\hand{54}}

% NOTE: Two small curved marks stand above and below the beginning of canceled Lorsque q=0. A long pencil stroke is in the left margin beside the final paragraph; no deletion is inferred from it.
